% ── SECTION 1: INTRODUCTION (~800 words) ─────────────────────────────────────

\section{Introduction: The Problem of Convergence}
\label{sec:intro}

The past century has produced two largely separate scholarly conversations.
In physics, interpretive engagement with quantum mechanics has produced
increasingly relational accounts of what physical reality is made
of~\citep{Rovelli2021,Bohm1980}.
In philosophy and religious studies, the relationship between scientific and
spiritual descriptions of reality has been debated~---~sometimes productively,
often at cross-purposes~\citep{Barbour2000,Polkinghorne1998}.

This paper argues that these conversations converge on a single finding,
arrived at independently by three traditions using three different methods:
\begin{quote}
\textit{Reality is not made of things. It is made of relationship. And the
force that constitutes and sustains those relationships --- at every scale
and in every domain --- is what each tradition, in its own vocabulary,
calls love.}
\end{quote}

The convergence claimed here is not analogical.
I am not suggesting that quantum entanglement metaphorically resembles
spiritual love, or that process philosophy poetically anticipates quantum
mechanics.
I claim structural isomorphism: the same logical architecture operating
in different vocabularies, arrived at by different methods, by traditions
that developed independently of one another.
The claim is stronger than analogy: each tradition's term for the relational
principle picks out the same structural feature of reality.

The argument employs what political philosophy calls overlapping
consensus~\citep{Rawls1993}: multiple independent theses, each valid
within its own domain, each pointing to the same underlying structure.
A physicist can validate Thesis~1 without accepting a word of philosophy
or spiritual writing.
A \bahai scholar can validate Thesis~3 without understanding density matrices.
The convergence itself --- the fact that all three arrive at the same structure
--- constitutes an argument more powerful than any single tradition can produce alone.

Two methodological clarifications follow.
The argument is not circular: it does not begin from the premise that
reality is relational and conclude that the force of relation is what
each tradition names.
It begins from three independent observations --- that the surviving
interpretations of quantum mechanics have a particular structure, that
Whitehead's process metaphysics has the same structure, and that the
\bahai writings have it as well --- each established within its own
tradition, by its own methods, prior to and without reference to the
others.
The convergence is then read off these three independently-grounded
observations; it is not built into them.
Nor is the argument vulnerable to the coincidence objection in the way
single-tradition arguments are.
Coincidence is a viable explanation for any one structural similarity;
coincidence as the explanation for six structural convergences
(\S\ref{sec:convergences}) arrived at across three civilizations and
roughly 170 years, using three epistemically distinct methods, becomes
progressively less plausible as the count rises.
The paper's weight rests not on the depth of any single convergence but
on the breadth of independent triangulation.

What gives this convergence its argumentative force is the chronological
sequence in which the three traditions arrived at their conclusions.
Table~\ref{tab:chronology} documents that sequence.

\begin{table}[h]
\centering
\caption{Chronological sequence of the three traditions' convergence.}
\label{tab:chronology}
\begin{tabularx}{\textwidth}{lXl}
\toprule
\textbf{Period} & \textbf{Development} & \textbf{Tradition} \\
\midrule
1850s--1892 & \Bahaullah articulates the relational constitution of existence
              and \mahabbat as the creative foundation of reality
            & \bahai \\
1892--1921  & \AbdulBaha develops the metaphysics of motion-as-life,
              relational affinity, and the non-composite soul
            & \bahai \\
1927        & Heisenberg uncertainty principle: no particle possesses
              simultaneously definite position and momentum
            & Physics \\
1929        & Whitehead, \textit{Process and Reality}: first systematic
              philosophical account of process ontology
            & Philosophy \\
1935        & Einstein-Podolsky-Rosen: formal challenge to the completeness
              of quantum mechanics raises questions about locality
            & Physics \\
1964        & Bell's theorem: formal proof that no local hidden-variable theory
              can reproduce quantum statistics
            & Physics \\
1982        & Aspect's experiments confirm Bell inequality violations;
              time-varying analyzers substantially address the locality loophole
            & Physics \\
1996        & Rovelli's relational quantum mechanics: quantum states are
              irreducibly observer-relative, never absolute
            & Physics \\
2004        & Epperson bridges Whitehead's process ontology and quantum mechanics,
              establishing their compatibility
            & Philosophy \\
2015        & First fully loophole-free Bell tests~\citep{Hensen2015,Giustina2015,Shalm2015}:
              locality and detection loopholes closed simultaneously
            & Physics \\
2022        & Aspect, Clauser, and Zeilinger share the Nobel Prize:
              recognition of the experimental confirmation program inaugurated
              by Bell's theorem
            & Physics \\
\bottomrule
\end{tabularx}
\end{table}

From the close of \Bahaullah's writings in 1892, the framework was named
philosophically by Whitehead within forty years and confirmed experimentally
via Bell's theorem within seventy.
This sequence does not prove supernatural authority --- that is not the
argument being made here.
What it establishes is that three independent methodologies arrived at
the same structural conclusion in a sequence that renders coincidence an
increasingly implausible explanation.

The paper proceeds as follows.
Section~\ref{sec:substance} demonstrates that quantum mechanics formally
requires process ontology, because Bell's theorem has eliminated
substance ontology as an internally consistent alternative.
Section~\ref{sec:whitehead} presents Whitehead's process ontology as the
philosophical framework that names the structure quantum mechanics requires.
Section~\ref{sec:bahai} demonstrates that the \bahai writings constitute
an independent prior instantiation of that same framework.
Section~\ref{sec:convergences} documents six structural isomorphisms across
all three traditions.
Section~\ref{sec:semantic} presents the semantic bridge: how the technical
terms of each tradition name the same underlying phenomenon.
Section~\ref{sec:rcf-phil} presents the Relational Coherence Framework (\RCF)
in philosophical language: the formal framework, developed in the companion
physics paper~\citep{Adams2026Physics}, that translates the convergence into
a quantitative object (the Affinity Tensor~$\alpha$), a testable two-qubit
gate ($\URA$), and experimentally discriminating predictions.
Section~\ref{sec:soul} explores the implications for the philosophy of
consciousness, with particular attention to the \bahai claim that the
soul is non-composite.
Section~\ref{sec:conclusion} draws out what the convergence means for
science-religion dialogue and the philosophy of science.
